\section{Diskussion}
\offen{Kapitel 7 -- Stichpunkte im Entwurf}
\offen{Eigener Agent als KI-System: Einstufung nach AI Act (vermutlich minimales Risiko, weil er Fachleuten nur Vorschläge macht)}
\offen{Datenhoheit: Nutzung eines US-Anbieters für das Sprachmodell, Bezug zu CLOUD Act und Art.~48 DSGVO (Abschnitt~2.2)}
\offen{Übertragbarkeit: Fall~4 mit anderer Rechtsordnung (Australien, keine übergeordnete Ebene wie die EU)}
\offen{Grenze: Das Organisationsprofil enthält nur Angaben, die für den Prozess wichtig sind. Ob der Agent mit einem größeren Profil voller unwichtiger Angaben genauso genau bleibt, wurde nicht geprüft (möglicher Zusatzversuch zu FF3)}
\offen{Grenze: Die Fachleute prüfen nacheinander, daher lässt sich keine Übereinstimmung zwischen unabhängigen Prüfungen (z.\,B. Cohens Kappa) angeben}
\offen{Grenze der Laufzeitkontrolle (Fall~2): Ein Monitor prüft, wann ein KI-System etwas tun darf (Regeln zum Ablauf, etwa Weitergabe der Planung erst nach Unterschrift der Fachfirma), aber nicht, ob sein Ergebnis fachlich stimmt (Regeln zum Inhalt, etwa Melderabstände nach DIN VDE 0833-2, die in einer nicht ladbaren Norm stehen und eine geometrische Auswertung des Grundrisses verlangen)}
